\section{Geometry of the ordered saddle}\label{sec:geometry}\label{paper:geometry}

The optimizer does more than select the exponential scale of the answer.
Its equality blocks determine which count windows are narrow, and its
first derivatives give the budgets available to the path constraints.
We begin with a general finite rank chain, since restrictions to equality
subspaces arise naturally when parameters move. We then specialize to the
complete chain needed for the divisor problem.

\subsection{The objective and the main conclusions}

Fix integers
\[
 H=h_1>h_2>\cdots>h_m\ge 2,
\]
where the one-element chain is allowed.  Set
\[
 p_i=h_i-h_{i+1}\quad(i<m),\qquad p_m=h_m-1,
 \qquad V_j=\sum_{i=1}^j L_i,\quad V_0=0,
\]
with arbitrary \(L_i>0\).  The symbols \(p_i\) always denote these
positive integers.  Child probabilities will instead be denoted by
\(\beta_i\).

On the ordered simplex
\[
 \calD=\{s\in\R^m:1\ge s_1\ge\cdots\ge s_m\ge0\}
\]
define, first at positive products,
\begin{align}
 G_m(s)&=s_m\log h_m,\nonumber\\
 G_i(s)&=s_i\log\bigl(p_i+\exp(G_{i+1}(s)/s_i)\bigr)
                   &&(i<m),\label{geom:eq:recursion}\\
 Z_i(s)&=\exp G_i(s),\nonumber\\
 J(L;s)&=\sum_iL_i+\sum_j s_jp_jV_j-\sum_iL_iZ_i(s).
                                                        \label{geom:eq:objective}
\end{align}
Use the continuous extension at zero products.  Indeed, induction gives
\(0\le G_i(s)\le s_i\log h_i\): if
\(G_{i+1}\le s_{i+1}\log h_{i+1}\), ordering gives
\(\exp(G_{i+1}/s_i)\le h_{i+1}\).  Thus when \(s_i=0\) its whole
suffix is zero and \(G_i=0\), \(Z_i=1\).  In particular,
\begin{equation}
 1\le Z_i(s)\le h_i \qquad(s\in\calD).             \label{geom:eq:globalZ}
\end{equation}

For \(E\subseteq\{1,\ldots,m-1\}\), let
\[
 \calD_E=\{s\in\calD:s_j=s_{j+1}\text{ for every }j\in E\}.
\]
Every maximum below is over this \emph{whole} feasible set.  Equalities
not specified by \(E\), and all other order inequalities, remain part
of the optimization.  Write \(s^E(L)\) for the maximizer and
\(J_E^*(L)=J(L;s^E(L))\).  Existence and uniqueness will be proved.

Throughout, \(c,C>0\) may change from line to line and depend only on
the fixed chain.  In particular they do not depend on \(L\), its ratios,
the face \(E\), or the sizes of product gaps and multipliers.  The notation
\(A\asymp B\) means \(cB\le A\le CB\), with this convention.

\begin{theorem}[Global bounds]\label{geom:thm:global}
Put
\begin{equation}
 c_0=2\log2-1,\qquad
 \varepsilon_H=\frac{c_0}{2H(\log H)^2}.             \label{geom:eq:epsilon}
\end{equation}
For every positive length vector and every \(E\), there is a unique
maximizer on \(\calD_E\), and it satisfies
\begin{equation}
 \frac12<s_m^E(L)\le s_1^E(L)<1-\varepsilon_H.
                                                        \label{geom:eq:outer}
\end{equation}
\end{theorem}

At a face maximizer define the signed cumulative multipliers
\begin{equation}
 \lambda_j^E(L)=-\sum_{i\le j}\partial_{s_i}J(L;s^E(L)),
 \qquad \lambda_0^E=\lambda_m^E=0.                 \label{geom:eq:lambda}
\end{equation}
The equality for \(\lambda_m^E\) follows from stationarity in the common
product direction.  For \(j\notin E\), these multipliers are nonnegative
and complementary to \(s_j^E-s_{j+1}^E\).  For \(j\in E\) they are signed.

\begin{theorem}[Stability across changing active faces]\label{geom:thm:stability}
Fix \(E\), let \(L'=L+\Delta\), and assume
\(|\Delta_i|\le L_i/2\).  Put \(D=\sum_i|\Delta_i|\).  Then
\begin{align}
 |s_i^E(L')-s_i^E(L)|
   &\le C\sum_{\ell=1}^m
                \frac{|\Delta_\ell|}{V_{\max(i,\ell)}},   \label{geom:eq:stability}\\
 |J_E^*(L')-J_E^*(L)|&\le CD,                           \label{geom:eq:value}\\
 |\lambda_j^E(L')-\lambda_j^E(L)|&\le CD.                \label{geom:eq:multiplier}
\end{align}
No extra lower threshold for the lengths is required.  The active
inequalities outside \(E\) may appear, disappear, or remain degenerate
along the perturbation.
\end{theorem}

For \(x\in\calD\), write \(\gap_j(x)=x_j-x_{j+1}\).
Let \(\bar s=s^E(L)\), \(s^*=s^\varnothing(L)\), and set
\begin{equation}
 a_j=
 \begin{cases}(-\lambda_j^E(L))_+,&j\in E,\\0,&j\notin E,\end{cases}
 \qquad \mathcal A_E=\sum_{j<m}\frac{a_j^2}{V_j}.       \label{geom:eq:AE}
\end{equation}

\begin{theorem}[Release of equalities]\label{geom:thm:opening}
Uniformly in \(L\) and \(E\),
\begin{align}
 J_\varnothing^*(L)-J_E^*(L)&\asymp\mathcal A_E,       \label{geom:eq:openingrate}\\
 \sum_i V_i(s_i^*-\bar s_i)^2&\asymp\mathcal A_E.     \label{geom:eq:openingdisp}
\end{align}
Both quantities are zero when \(\mathcal A_E=0\).
For \(E=\{j\}\), put \(a=(-\lambda_j^E)_+\).  If \(a=0\),
then \(\bar s=s^*\).  If \(a>0\), additionally
\begin{equation}
 \gap_j(s^*)\asymp \frac a{V_j},\qquad
 J_\varnothing^*(L)-J_{\{j\}}^*(L)\asymp\frac{a^2}{V_j}.
                                                        \label{geom:eq:single}
\end{equation}
The individual-gap assertion is only made for a single released edge.
\end{theorem}

We prove these estimates directly from the finite objective. The curvature
controls each suffix on its own cumulative length scale; this is the
reason that perturbations remain controlled when the active equality
constraints change.

\subsection{Concavity and the outer boundaries}

The map \(x\mapsto\log(p+e^x)\) is convex and increasing.  Its
perspective \((s,y)\mapsto s\log(p+e^{y/s})\) is convex for \(s>0\)
and is increasing in \(y\).  Induction in \eqref{geom:eq:recursion} therefore
shows that each \(G_i\), and hence each \(Z_i\), is convex.  The continuous
extension gives convexity on \(\calD\).  Thus \(J\) is concave and
continuous on the compact convex set \(\calD_E\), so a maximum exists.

\subsubsection{The bottom plateau}

\begin{lemma}\label{geom:lem:bottom}
Every maximizer on \(\calD_E\) satisfies \(s_m>1/2\).
\end{lemma}

\begin{proof}
Suppose that \(a,\ldots,m\) is the maximal bottom block of equal
products, with common value \(\sigma\le1/2\).  Increasing this entire
block by a small positive amount is feasible.  If \(a>1\), its preceding
gap is strict, so no prescribed equality crosses that gap.  Exact
telescoping gives
\[
 G_i=\sigma\log h_i,\qquad Z_i=h_i^\sigma\quad(i\ge a).
\]
For \(i<a\), set
\[
 \beta_i=\frac{\exp(G_{i+1}/s_i)}
                    {p_i+\exp(G_{i+1}/s_i)}.
\]
Along this bottom-block direction,
\[
 \frac{dG_i}{d\sigma}
   =\left(\prod_{i\le b<a}\beta_b\right)\log h_a.
\]
The identity
\[
 Z_i\beta_i=Z_{i+1}\beta_i^{\,1-s_i}
\]
implies
\[
 Z_i\prod_{i\le b<a}\beta_b
   = Z_a\prod_{i\le b<a}\beta_b^{\,1-s_b}\le Z_a.
\]
When \(\sigma=0\), all preceding products, if present, are positive,
and these formulas give the right derivative.  There are no preceding
terms if \(a=1\).

The derivative of the linear budget is
\[
 \sum_{j\ge a}p_jV_j
   =(h_a-1)V_{a-1}+\sum_{i\ge a}(h_i-1)L_i.
\]
Consequently the directional derivative of \(J\) is at least
\begin{equation}
 V_{a-1}(h_a-1-h_a^\sigma\log h_a)
       +\sum_{i\ge a}L_i(h_i-1-h_i^\sigma\log h_i). \label{geom:eq:bottomderivative}
\end{equation}
For every real \(h>1\), one has \(\sqrt h\log h<h-1\).
Indeed, with \(t=\sqrt h\), the difference
\(t^2-1-2t\log t\) vanishes at \(1\) and has positive derivative
\(2(t-1-\log t)\) for \(t>1\).
Every bracket in \eqref{geom:eq:bottomderivative} is thus positive.
Since the lengths are positive, the derivative is strictly positive.
This contradicts maximality, including at \(\sigma=0\).
\end{proof}

\subsubsection{The top plateau}

\begin{lemma}\label{geom:lem:top}
Every maximizer on \(\calD_E\) satisfies \(s_1<1-\varepsilon_H\).
\end{lemma}

\begin{proof}
Let \(1,\ldots,b\) be the maximal top equal-product block, with
common value \(\sigma\).  If \(b<m\), put
\(C=h_{b+1}\) and \(z=\exp(G_{b+1}/\sigma)\); otherwise put \(C=z=1\).
For \(i\le b\), let \(n_i=h_i-C\ge1\).  Telescoping yields
\[
 G_i=\sigma\log(n_i+z).
\]
Here \(1\le z\le C\) and \(n_i+z\le H\), by the ordered recursion.
Varying the top block while keeping its child products fixed gives
\begin{equation}
 \frac{dZ_i}{d\sigma}=(n_i+z)^\sigma A(n_i,z),\qquad
 A(n,z)=\log(n+z)-\frac{z\log z}{n+z}.             \label{geom:eq:topderivative}
\end{equation}
This formula includes the variation of \(z=\exp(G_{b+1}/\sigma)\).
The budget derivative is \(\sum_{i\le b}L_i n_i\).

Write
\[
 E_n(z)=(n+z)\log(n+z)-z\log z=(n+z)A(n,z).
\]
It is increasing in \(z\), since \(\partial_zE_n=\log(1+n/z)>0\).
Moreover
\[
 E_n(z)\ge E_n(1)=(n+1)\log(n+1)\ge2n\log2.
\]
For the last inequality,
\((n+1)\log(n+1)/n\) is increasing for \(n>0\), with derivative
\([n-\log(n+1)]/n^2>0\), and \(n\ge1\).
Also \(0<A(n,z)\le\log H\).  With \(n,z\) fixed at their values
at the point under consideration, for \(0\le\tau\le1\),
\[
 \frac{d}{d\tau}\bigl((n+z)^\tau A(n,z)\bigr)
 \le H(\log H)^2.
\]
This auxiliary comparison keeps \(z\) fixed.
If \(\sigma\ge1-\varepsilon_H\), it follows that
\[
 (n+z)^\sigma A(n,z)-n
 \ge n(2\log2-1)-\varepsilon_HH(\log H)^2
 \ge c_0/2>0.
\]
Thus raising the top block has strictly negative derivative for \(J\).
Lowering it a sufficiently small amount is feasible and increases
\(J\), a contradiction.  This also covers \(\sigma=1\).
\end{proof}

Let \(\varepsilon=\varepsilon_H\), which lies in \((0,1/2)\), and put
\[
 \calD_\varepsilon
   =\{s:\varepsilon\le s_m\le\cdots\le s_1\le1-\varepsilon\}.
\]
The preceding lemmas put every maximizer on every prescribed face in
this single compact set.  We will use differential estimates on the
slightly larger convex set \(\calD_{\varepsilon/2}\).

\subsection{Full suffix curvature and uniqueness}

For \(i<m\), introduce
\[
 a_i=G_i/s_i=\log S_i,\qquad
 u_i=s_{i+1}/s_i,\qquad
 \beta_i=\frac{e^{G_{i+1}/s_i}}{S_i},\qquad
 A_i=a_i-u_i\beta_i a_{i+1},
\]
where \(S_i=p_i+e^{G_{i+1}/s_i}\); set
\(a_m=A_m=\log h_m\).
All these functions and the derivatives needed below are bounded on
\(\calD_{\varepsilon/2}\).  The quantities \(s_i,a_i,\beta_i,1-\beta_i\)
have positive chain-dependent lower bounds.  The same holds for \(A_i\):
it is the entropy of the distribution consisting of \(p_i\) named
probabilities \(1/S_i\) and the one compound probability \(\beta_i\).
These assertions use compactness but no separation of adjacent products.

\begin{lemma}[One perspective term controls its entire suffix]
\label{geom:lem:suffix}
For any \(s\in\calD_{\varepsilon/2}\), direction \(z\in\R^m\), and \(i\),
\begin{equation}
 c\sum_{j\ge i}z_j^2\le D^2Z_i(s)[z,z]
                   \le C\sum_{j\ge i}z_j^2.      \label{geom:eq:suffix}
\end{equation}
\end{lemma}

\begin{proof}
Let \(g_i=DG_i[z]\) and \(Q_i=D^2G_i[z,z]\).  Differentiating the
perspective recursion gives
\begin{align}
 g_m&=(\log h_m)z_m,\qquad Q_m=0,\nonumber\\
 g_i&=A_i z_i+\beta_i g_{i+1},\nonumber\\
 Q_i&=\beta_iQ_{i+1}
  +\frac{\beta_i(1-\beta_i)}{s_i}
              (g_{i+1}-u_i a_{i+1}z_i)^2.        \label{geom:eq:perspectivecurvature}
\end{align}
In particular \(Q_i\ge0\).  Put \(E_i=g_i^2+Q_i\) and
\(w_i=g_{i+1}-u_i a_{i+1}z_i\).  Then
\[
 g_i=a_i z_i+\beta_iw_i,\qquad
 Q_i\ge c w_i^2+cQ_{i+1}.
\]
Thus \(z_i^2\le CE_i\), and then \(g_{i+1}^2\le CE_i\), so that
\(E_{i+1}\le CE_i\).  Iteration over the fixed number of indices,
starting with \(E_m=(\log h_m)^2z_m^2\), proves
\[
 E_i\ge c\sum_{j\ge i}z_j^2.
\]
The reverse inequality follows recursively from
\eqref{geom:eq:perspectivecurvature} and its bounded coefficients.
Finally,
\[
 D^2Z_i[z,z]=Z_iE_i,
\]
and \(Z_i\) is bounded above and below.
\end{proof}

Define the negative objective Hessian by
\(\calH(L;s)=-D_s^2J(L;s)\).
Summing \eqref{geom:eq:suffix} with weights \(L_i\) gives
\begin{equation}
 c\sum_jV_j z_j^2
 \le z^{\mathsf T}\calH(L;s)z
 \le C\sum_jV_j z_j^2.                            \label{geom:eq:hessian}
\end{equation}
Furthermore,
\begin{equation}
 |\calH_{ij}(L;s)|\le C V_{\min(i,j)},             \label{geom:eq:hessianentry}
\end{equation}
because only terms \(L_\ell Z_\ell\) with \(\ell\le\min(i,j)\)
depend on both coordinates, and their second derivatives are bounded.
These estimates hold for every positive length vector, without
stationarity or a length hierarchy.

If two distinct maximizers existed on \(\calD_E\), concavity would
make their whole connecting segment maximizing.  The global boundary
lemmas put that segment in \(\calD_\varepsilon\), whereas
\eqref{geom:eq:hessian} makes the restriction to the segment strictly concave.
This contradiction proves uniqueness and finishes
Theorem~\ref{geom:thm:global}.

\subsection{An LDL inverse estimate on every equality subspace}

\begin{lemma}\label{geom:lem:inverse}
Suppose that a symmetric positive definite \(r\times r\) matrix \(A\)
and increasing positive numbers \(W_1,\ldots,W_r\) satisfy
\[
 c\diag(W)\le A\le C\diag(W),\qquad
 |A_{ij}|\le C W_{\min(i,j)}.
\]
Then
\begin{equation}
 |(A^{-1})_{ij}|\le\frac{C'}{W_{\max(i,j)}},       \label{geom:eq:inverse}
\end{equation}
where \(C'\) depends only on \(r,c,C\).
\end{lemma}

\begin{proof}
Perform scalar LDL elimination in the order \(1,\ldots,r\).
Every Schur complement retains the same diagonal quadratic bounds:
its quadratic form is obtained by minimizing the preceding form over
the eliminated coordinates.  The lower bound follows by discarding
the nonnegative eliminated diagonal terms; for the upper bound choose
those coordinates to be zero.  Thus the pivot at index \(i\) is
comparable to \(W_i\).

The entry estimate also persists, with constants depending on the
number of steps.  When index \(i\) is eliminated, the update in entry
\(j,\ell>i\) is bounded by
\[
 C\,\frac{W_iW_i}{W_i}=CW_i
       \le CW_{\min(j,\ell)}.
\]
Every multiplier in the unit lower triangular elimination matrix is
therefore bounded.  Its inverse has bounded entries as well, by the
finite nilpotent series for the strictly lower triangular part.
If \(A=\mathsf L\,\mathsf D\,\mathsf L^{\mathsf T}\), the \(i,j\) entry
of \(A^{-1}=\mathsf L^{-\mathsf T}\mathsf D^{-1}\mathsf L^{-1}\)
is a sum over indices \(t\ge\max(i,j)\) of bounded numerators divided
by pivots comparable to \(W_t\).  Monotonicity of \(W_t\) proves
\eqref{geom:eq:inverse}.
\end{proof}

Partition \(\{1,\ldots,m\}\) into consecutive blocks \(I_a=[b_a,e_a]\).
Let \(B\) map block coordinates to vectors constant on each block.
The restricted Hessian is \(A=B^{\mathsf T}\calH B\).
By \eqref{geom:eq:hessian}, its natural quadratic weights are
\(\sum_{j\in I_a}V_j\), which lie between \(V_{e_a}\) and \(mV_{e_a}\).
For \(a<b\), summing \eqref{geom:eq:hessianentry} over the two blocks gives
\(|A_{ab}|\le C V_{e_a}\).  Lemma~\ref{geom:lem:inverse} consequently gives
\begin{equation}
 \left|\bigl((B^{\mathsf T}\calH B)^{-1}\bigr)_{ab}\right|
       \le \frac C{V_{\max(e_a,e_b)}}.            \label{geom:eq:restrictedinverse}
\end{equation}
There are only finitely many such partitions, so a single constant
works for all of them.  The Hessian is that of the original objective
restricted to its current equality subspace.

\subsection{Moving optimizers and differentiation across faces}

We first record the elementary KKT relations with the sign convention
in \eqref{geom:eq:lambda}.  At a maximizer on \(\calD_E\), a small shift of
all products in either direction is feasible by
Theorem~\ref{geom:thm:global}; hence \(\sum_i\partial_iJ=0\).
For \(j\notin E\), increasing the first \(j\) coordinates is feasible
even when \(\gap_j=0\).  Its derivative is \(-\lambda_j\), so
\(\lambda_j\ge0\).  If \(\gap_j>0\), shifting that prefix in either
direction is feasible, giving \(\lambda_j=0\).  Thus
\begin{equation}
 \partial_iJ=\lambda_{i-1}-\lambda_i,\qquad
 \lambda_j\ge0,\quad\lambda_j\gap_j=0\quad(j\notin E).
                                                        \label{geom:eq:KKT}
\end{equation}
The multipliers at edges in \(E\) are not sign constrained.
Conversely, \eqref{geom:eq:KKT}, the signed equalities, and concavity
give the corresponding global variational inequality.

Fix \(E\), and consider \(L(t)=L+t\Delta\), \(0\le t\le1\), with
\(|\Delta_i|\le L_i/2\).  Then
\begin{equation}
 \tfrac12 V_j\le V_j(t)\le\tfrac32 V_j.            \label{geom:eq:Vcompare}
\end{equation}
Let \(s(t)=s^E(L(t))\).  All of these true constrained maxima lie in
\(\calD_\varepsilon\), by the global theorem.

\begin{lemma}\label{geom:lem:AC}
The path \(s(t)\), and all gradient components along it, are Lipschitz
and hence absolutely continuous.  At almost every \(t\), its derivative
is constant on each maximal block of equal current products.
\end{lemma}

\begin{proof}
Continuity follows from compactness and uniqueness: any convergent
subsequence of maximizing points as \(t_n\to t\) maximizes the limiting
objective, so its limit is \(s(t)\).

For \(t,t'\), the two constrained variational inequalities and strong
concavity on the segment joining the maximizers give, with
\(d=s(t')-s(t)\),
\[
 cV_1\|d\|^2
 \le \left|\bigl(\nabla J(L(t');s(t'))
                 -\nabla J(L(t);s(t'))\bigr)\cdot d\right|
 \le C|t'-t|D\|d\|.
\]
Here the first bound uses \eqref{geom:eq:hessian} and
\eqref{geom:eq:Vcompare}, and the second uses bounded length derivatives
of the gradient.  This Lipschitz bound establishes the regularity
needed for the componentwise estimate below.
The gradient path is also Lipschitz since the objective is smooth
on the compact product set and the length path is affine.

A nonnegative absolutely continuous function has derivative zero
at almost every point of its zero set.  To see this, at a density
point of the zero set where its derivative exists, take difference
quotients along that set.  The density theorem and almost-everywhere
differentiability remove the remaining null sets.
Apply this observation to each gap \(s_i(t)-s_{i+1}(t)\).
There are finitely many gaps, so at almost every common differentiability
point, \(\dot s(t)\) is constant on every current equal-product block.
\end{proof}

At such a time, let \(I_a=[b_a,e_a]\) be the current maximal blocks,
and write \(\dot s=B\dot q\).  Both exterior gaps of a block, when
present, are strictly positive.  They stay positive in a neighborhood
of this time, and their multipliers vanish throughout that neighborhood.
No prescribed equality crosses one of these strict gaps.  Consequently,
keeping these block endpoints fixed, one has in the whole neighborhood
\begin{equation}
 \sum_{i\in I_a}\partial_iJ(L(t);s(t))=0.          \label{geom:eq:blockstationarity}
\end{equation}
This identity remains true if internal equalities split nearby.
It may therefore be differentiated at the chosen almost-everywhere
point.  This is the step that avoids assuming a differentiable
active-set chart.

Since \(J\) is linear in \(L\), define
\begin{equation}
 U_\ell(s)=\partial_{L_\ell}J
    =1+\sum_{j\ge\ell}s_jp_j-Z_\ell(s),\qquad
 b=\sum_\ell\Delta_\ell\nabla U_\ell(s(t)).
                                                        \label{geom:eq:U}
\end{equation}
Differentiating \eqref{geom:eq:blockstationarity} gives
\begin{equation}
 (B^{\mathsf T}\calH B)\dot q=B^{\mathsf T}b.     \label{geom:eq:blockODE}
\end{equation}
The vector \(\nabla U_\ell\) has zero entries before \(\ell\) and
uniformly bounded entries thereafter.  Its sum over a block can
therefore be nonzero only when that block's right endpoint is at
least \(\ell\).  By \eqref{geom:eq:restrictedinverse}, for \(i\in I_a\),
\begin{align}
 |\dot s_i(t)|
 &\le C\sum_\ell|\Delta_\ell|
      \sum_{\substack{b:\ e_b\ge\ell}}
                   \frac1{V_{\max(e_a,e_b)}(t)}\nonumber\\
 &\le C\sum_\ell\frac{|\Delta_\ell|}
                         {V_{\max(i,\ell)}(t)}.  \label{geom:eq:derivstability}
\end{align}
Indeed \(e_a\ge i\), \(e_b\ge\ell\), and there are at most \(m\)
blocks.  Integrating this almost-everywhere estimate and using
\eqref{geom:eq:Vcompare} proves \eqref{geom:eq:stability}.

\begin{remark}
An equality with zero multiplier may persist on a set of positive
measure, split, or reform.  The preceding argument includes all such
possibilities: tangency comes from absolute continuity of the gaps,
and the differentiated stationarity equation uses only the current
block's strictly separated exterior boundaries.
\end{remark}

\subsection{Values, multipliers, and own-scale parameters}

The coefficients \(U_\ell\) in \eqref{geom:eq:U} are bounded on the entire
original simplex, using \eqref{geom:eq:globalZ}.  Hence
\[
 |J(L';s)-J(L;s)|\le CD\qquad(s\in\calD).
\]
Taking the maximum on the same \(\calD_E\) in both directions proves
\eqref{geom:eq:value}.  In fact this value bound does not require the
relative perturbation assumption.

For the signed multiplier, absolute continuity and
\eqref{geom:eq:U} give almost everywhere
\[
 \dot\lambda_j
   =-\sum_{i\le j}b_i
       +\sum_{i\le j}\sum_a\calH_{ia}\dot s_a.
\]
The first term is \(O(D)\).  By \eqref{geom:eq:hessianentry} and
\eqref{geom:eq:derivstability}, the second term is at most
\[
 C\sum_{i\le j}\sum_\ell|\Delta_\ell|
          \sum_a\frac{V_{\min(i,a)}(t)}
                       {V_{\max(a,\ell)}(t)}.
\]
Every fraction is at most one.  The finite sums are therefore
\(O(D)\), uniformly in all length ratios.  Integration proves
\eqref{geom:eq:multiplier} and finishes Theorem~\ref{geom:thm:stability}.
In particular the bound applies to the unrestricted signed
multipliers at prescribed equality edges.

\begin{corollary}[Own-scale exponent stability]\label{geom:cor:ownscale}
Let \(F_i(s_i,\ldots,s_m)\) be bounded and Lipschitz on
\(\calD_{\varepsilon/2}\), with chain-dependent bounds.  Under the
hypotheses of Theorem~\ref{geom:thm:stability},
\begin{equation}
 |F_i(s^E(L'))-F_i(s^E(L))|\le \frac{CD}{V_i}.
                                                        \label{geom:eq:ownscale}
\end{equation}
For \(D\) bounded, each of the factors
\[
 (1+L_i)^{F_i(s^E(L))},\qquad
 \left(\frac{1+V_{i-1}}{1+L_i}\right)^{F_i(s^E(L))}
\]
changes by a bounded multiplicative factor when \(L\) is replaced
by \(L'\).  The bound does not involve a logarithm of a more remote
length.
\end{corollary}

\begin{proof}
For every suffix coordinate \(j\ge i\), \eqref{geom:eq:stability} is
bounded above by \(CD/V_i\), proving \eqref{geom:eq:ownscale}.
Also
\[
 \log(1+L_i)\le\log(1+V_i),\qquad
 \left|\log\frac{1+V_{i-1}}{1+L_i}\right|
      \le2\log(1+V_i).
\]
Multiply these bounds by \eqref{geom:eq:ownscale} and use
\(\log(1+V_i)/V_i\le1\).
The bases themselves change by bounded multiplicative factors,
because their numerator and denominator change by bounded additive
amounts and each has a \(1+\) normalization.  Boundedness of \(F_i\)
finishes the proof.
\end{proof}

For example the ratio parameters \(d_i=s_i/s_{i+1}-1\) satisfy
\[
 |d_i(L')-d_i(L)|\le CD/V_i,
\]
since products stay away from zero.  This observation and
\eqref{geom:eq:multiplier} are separate statements.  They do not by
themselves justify stability of a proposed persistence expression
involving division by a shorter native length or a change of formula
between different platforms.

\subsection{Opening a single ordered-product edge}

Fix \(1\le j<m\).  Let \(\bar s\) maximize with the sole extra equality
\(s_j=s_{j+1}\), and let \(s^*\) be the unrestricted ordered maximizer.
All other inequalities are optimized in both problems.
Write \(\bar\lambda_i=-\sum_{a\le i}\partial_aJ(L;\bar s)\) and
\(a=(-\bar\lambda_j)_+\).

If \(\bar\lambda_j\ge0\), every multiplier has the sign and
complementarity required for the unrestricted problem.  More explicitly,
for every \(y\in\calD\), summation by parts gives
\[
 \nabla J(L;\bar s)\cdot(y-\bar s)
   =-\sum_{i<m}\bar\lambda_i
                  \bigl(\gap_i(y)-\gap_i(\bar s)\bigr)\le0.
\]
Concavity and uniqueness imply \(\bar s=s^*\).

Suppose now that \(a>0\).  Let \(e^{(j)}\) have its first \(j\)
coordinates equal to one and the rest zero.  The perturbation
\(\bar s+\tau e^{(j)}\) opens only gap \(j\); every other gap is
unchanged.  It is feasible for all sufficiently small \(\tau>0\),
even if some preceding or following inequalities are equalities.
Its first derivative at zero is
\[
 \nabla J(L;\bar s)\cdot e^{(j)}=-\bar\lambda_j=a.
\]
On \(\calD_{\varepsilon/2}\), its negative second derivative is at most
\[
 C\sum_{i\le j}V_i\le CmV_j.
\]
We also have the useful elementary bound
\begin{equation}
 |\bar\lambda_j|\le C_0V_j.                       \label{geom:eq:lambdabound}
\end{equation}
Indeed
\(\partial_iJ=p_iV_i-\sum_{\ell\le i}L_\ell\partial_iZ_\ell\)
has absolute value at most \(CV_i\), by the compact derivative
bounds, and there are at most \(m\) terms in \(\bar\lambda_j\).

Choose a sufficiently large chain-dependent \(C_1\) and set
\(\tau=a/(C_1V_j)\).  By \eqref{geom:eq:lambdabound}, this can be made
at most \(\varepsilon/2\), so the ray stays in
\(\calD_{\varepsilon/2}\).  Taylor's formula with integral remainder
then gives
\begin{equation}
 J(L;\bar s+\tau e^{(j)})-J(L;\bar s)
    \ge a\tau-CV_j\tau^2\ge c\,\frac{a^2}{V_j}.
                                                        \label{geom:eq:singlelower}
\end{equation}

For the matching bound, put \(\Delta s=s^*-\bar s\) and
\(g=\gap_j(s^*)\).  Since the original gap is zero,
\(g=\Delta s_j-\Delta s_{j+1}\ge0\).  Summation by parts and
complementarity at all other edges yield
\begin{align}
 \nabla J(L;\bar s)\cdot\Delta s
  &=-\sum_{i<m}\bar\lambda_i
                  \bigl(\gap_i(s^*)-\gap_i(\bar s)\bigr)\nonumber\\
  &=ag-\sum_{i\ne j}\bar\lambda_i\gap_i(s^*)\le ag.
                                                        \label{geom:eq:singlelinear}
\end{align}
The neighboring nonnegative multipliers retain their negative
contribution.
Moreover
\[
 \sum_iV_i(\Delta s_i)^2
 \ge V_j\bigl((\Delta s_j)^2+(\Delta s_{j+1})^2\bigr)
 \ge \tfrac12V_jg^2.
\]
Strong concavity integrated along the segment between the two
maximizers therefore implies
\begin{equation}
 0\le J(L;s^*)-J(L;\bar s)
       \le ag-cV_jg^2\le C\,\frac{a^2}{V_j}.      \label{geom:eq:singleupper}
\end{equation}
Combining \eqref{geom:eq:singlelower} with the bound by \(ag\) gives
\(g\ge ca/V_j\).  Nonnegativity in \eqref{geom:eq:singleupper} gives the
reverse inequality.  This proves \eqref{geom:eq:single}.

The natural bounded-rate opening coordinate is thus
\[
 \frac{a}{\sqrt{V_j}}.
\]
This is a statement about the exact optimized rate.  If this coordinate
diverges, using the restricted-face rate instead of the unrestricted
one produces an unbounded exponential discrepancy.

\subsection{Simultaneous release of a prescribed face}

We now allow an arbitrary \(E\), with the notation
\eqref{geom:eq:AE}.  All signed multipliers refer to the maximum
\(\bar s=s^E(L)\) of the original objective on that face.

\subsubsection{A feasible test direction}

For each \(j\) with \(a_j>0\), choose
\[
 \tau_j=\frac{a_j}{C_1V_j},\qquad
 z=\sum_{j:a_j>0}\tau_j e^{(j)}.
\]
This opens exactly these gaps, because
\(z_j-z_{j+1}=\tau_j\), with \(\tau_j=0\) at all other edges.
It preserves the order constraints.  The argument for
\eqref{geom:eq:lambdabound} gives \(a_j\le C_0V_j\), so
\(\sum_j\tau_j\le mC_0/C_1\).  Choose \(C_1\) large enough that
the entire test ray stays in \(\calD_{\varepsilon/2}\).

Its linear gain is exactly
\[
 \nabla J(L;\bar s)\cdot z=\sum_j a_j\tau_j.
\]
For the quadratic cost, Cauchy--Schwarz in the finite prefix sum gives
\begin{align}
 \sum_iV_i z_i^2
 &\le m\sum_j\tau_j^2\sum_{i\le j}V_i
 \le m^2\sum_jV_j\tau_j^2.                       \label{geom:eq:multiquad}
\end{align}
Taylor's formula and \eqref{geom:eq:hessian}, after one further
chain-dependent enlargement of \(C_1\), yield
\begin{equation}
 J(L;\bar s+z)-J(L;\bar s)\ge c\mathcal A_E.      \label{geom:eq:multilower}
\end{equation}
The unrestricted gain is therefore at least \(c\mathcal A_E\).

\subsubsection{The upper bound and the total displacement}

Put \(\Delta s=s^*-\bar s\) and
\(Q=\sum_iV_i(\Delta s_i)^2\).
At edges outside \(E\) the multipliers are nonnegative and
complementary; at edges in \(E\) the original gap is zero.
Summation by parts consequently gives
\begin{equation}
 \nabla J(L;\bar s)\cdot\Delta s
 \le\sum_{j:a_j>0}a_j\gap_j(s^*).                 \label{geom:eq:multilinear}
\end{equation}
For \(j\in E\), \(\gap_j(s^*)=\Delta s_j-\Delta s_{j+1}\), so
\begin{equation}
 \sum_{j\in E}V_j\gap_j(s^*)^2
 \le2\sum_{j<m}V_j\bigl((\Delta s_j)^2+
                              (\Delta s_{j+1})^2\bigr)
 \le4Q.                                         \label{geom:eq:gapQ}
\end{equation}
The last bound only uses \(V_j\le V_{j+1}\).
Weighted Cauchy--Schwarz in \eqref{geom:eq:multilinear}, followed by
integrated strong concavity, proves
\begin{equation}
 0\le J(L;s^*)-J(L;\bar s)
             \le2\sqrt{\mathcal A_EQ}-cQ.       \label{geom:eq:multiupper}
\end{equation}
Thus \(Q\le C\mathcal A_E\), and the rate gain is at most
\(C\mathcal A_E\).  Conversely, \eqref{geom:eq:multilower} and
\eqref{geom:eq:multiupper} with its negative term omitted give
\(Q\ge c\mathcal A_E\) when \(\mathcal A_E>0\).
If \(\mathcal A_E=0\), \eqref{geom:eq:multiupper} forces \(Q=0\),
and both optima coincide.  This proves
\eqref{geom:eq:openingrate}--\eqref{geom:eq:openingdisp} and completes
Theorem~\ref{geom:thm:opening}.

\begin{remark}
The simultaneous comparison gives a weighted total displacement,
but does not assert
\(\gap_j(s^*)\asymp a_j/V_j\) for every member of a multi-edge release.
Several signed multipliers may interact through the order constraints.
\end{remark}


\subsection{The finite optimization and its product bounds}

Fix integers
\[
 H=h_1>h_2>\cdots>h_m\ge2,
 \qquad p_j=h_j-h_{j+1}\ (j<m),\qquad p_m=h_m-1.
\]
Let $L_i>0$, $V_j=\sum_{i\le j}L_i$, and $V_0=0$. The product
variables belong to the ordered simplex
\[
 \mathcal S=\{1\ge s_1\ge s_2\ge\cdots\ge s_m\ge0\}.
\]
For positive products define recursively
\begin{align}
 G_m&=s_m\log h_m,&
 G_j&=s_j\log\bigl(p_j+\exp(G_{j+1}/s_j)\bigr)\quad(j<m),
 \label{hier:eq:G}\\
 Z_j&=e^{G_j},&
 J(\mathbf s)&=\sum_iL_i+\sum_j s_jp_jV_j-\sum_iL_iZ_i.
 \label{hier:eq:J}
\end{align}
The functions extend continuously to $\mathcal S$: if $s_j=0$, the
entire suffix has zero product and we set $G_j=0$. Indeed, induction
gives $0\le G_j\le s_j\log h_j$, which proves continuity at those
faces. A maximizer exists by compactness. The perspective of
$y\mapsto\log(p+e^y)$ is convex and nondecreasing in its second
argument; recursive perspective composition shows that each $G_j$,
then each $Z_j$, is convex on $\mathcal S$. Thus $J$ is concave. We
use the unique global maximizer furnished by Theorem~\ref{geom:thm:global}.

When products are positive, put
\begin{align}
 S_m&=h_m,& S_j&=p_j+S_{j+1}^{u_j},&u_j&=s_{j+1}/s_j,\\
 \beta_j&=S_{j+1}^{u_j}/S_j,&
 A_j&=\log S_j-u_j\beta_j\log S_{j+1}\quad(j<m),&
 A_m&=\log h_m. \label{hier:eq:node}
\end{align}
Then $G_j=s_j\log S_j$, $Z_j=S_j^{s_j}$, and
$2\le S_j\le h_j\le H$. We repeatedly use the entropy increment
\begin{equation}
 E_n(z)=(n+z)\log(n+z)-z\log z,
 \qquad n>0,\ z\ge1. \label{hier:eq:entropy}
\end{equation}
For a nonterminal node, $A_j=E_{p_j}(S_{j+1}^{u_j})/S_j>0$.

\begin{lemma}[Uniform upper product bound]\label{hier:lem:upper}
Every maximizer satisfies $s_1<1-\epsilon_H$, where
$\epsilon_H=(2\log2-1)/(2H(\log H)^2)$.
\end{lemma}
\begin{proof}
This is the upper bound in Theorem~\ref{geom:thm:global}, applied with
no prescribed equalities.
\end{proof}

\begin{lemma}[Canonical thresholds]\label{hier:lem:alpha}
For real $h>1$ define
\[
 \alpha_h=\frac{\log((h-1)/\log h)}{\log h}
 =1-\frac{\log(h\log h/(h-1))}{\log h}.
\]
Then $\alpha_h>1/2$ and $\alpha_h$ is strictly increasing.
In particular, $2^{\alpha_2}\log2=1$.
\end{lemma}
\begin{proof}
Write $y=\log h$ and $z=y/2$. Direct calculation gives
\[
 \alpha_h=\frac12+\frac{\log(\sinh z/z)}{2z}.
\]
The second term is positive. Its derivative has the sign of
\[
 N(z)=z\coth z-1-\log(\sinh z/z).
\]
Here $N(0+)=0$ and $N'(z)=1/z-z/\sinh^2z>0$. This proves both
assertions without a discrete approximation.
\end{proof}

\begin{lemma}[Transported natural intensities]\label{hier:lem:transport}
For $i<j$ at positive products,
\begin{equation}
 Z_i\prod_{q=i}^{j-1}\beta_q
 =Z_j\prod_{q=i}^{j-1}\beta_q^{\,1-s_q}\le Z_j.
 \label{hier:eq:transport}
\end{equation}
If $s_1<1$, the inequality is strict.
\end{lemma}
\begin{proof}
The one-edge identity $Z_q\beta_q=Z_{q+1}\beta_q^{1-s_q}$ follows
directly from \eqref{hier:eq:G}--\eqref{hier:eq:node}. Multiplication telescopes.
Since $0<\beta_q<1$, the conclusions follow.
\end{proof}

\begin{theorem}[Global product compactness]\label{hier:thm:compact}
For $H\ge3$ and every positive length vector, every maximizer satisfies
\begin{equation}
 \alpha_2<s_m\le s_1<1-\epsilon_H. \label{hier:eq:compact}
\end{equation}
For the exceptional one-node chain $H=2$, the unique maximizing
product is $\alpha_2$. In particular all products, local ratios and
finite positive mark probabilities used below lie in compact sets
determined by $H$.
\end{theorem}
\begin{proof}
Lemma~\ref{hier:lem:upper} already gives the upper bound. Let $a,\ldots,m$
be the bottom maximal platform, with product $\sigma\le\alpha_2$.
Raising its common product is feasible. Put
\[
 I_{\rm in}=\sum_{i<a}L_iZ_i\prod_{q=i}^{a-1}\beta_q.
\]
The derivative in that direction is
\begin{equation}
 (h_a-1)V_{a-1}-I_{\rm in}\log h_a
 +\sum_{i\ge a}L_i\bigl(h_i-1-h_i^\sigma\log h_i\bigr).
 \label{hier:eq:bottomderiv}
\end{equation}
This also holds as a one-sided derivative at $\sigma=0$: the older
products are positive, and the bottom recursion is exactly
$G_i=\sigma\log h_i$. Transport gives
$I_{\rm in}\le h_a^\sigma V_{a-1}$. By
Lemma~\ref{hier:lem:alpha}, all brackets in \eqref{hier:eq:bottomderiv} are
nonnegative, and a bracket with $h_i>2$ is strictly positive.
Thus \eqref{hier:eq:bottomderiv} is positive if such a rank occurs in
the bottom platform. Otherwise the platform is the single rank
$h_m=2$, with $a=m>1$ when $H\ge3$. In that case
$I_{\rm in}<2^\sigma V_{m-1}$ by Lemma~\ref{hier:lem:upper} and the
strict transport identity, so the derivative is again positive,
even at $\sigma=\alpha_2$. Both alternatives contradict maximality.
For the chain $H=2$, directly differentiating
$L(1+s-2^s)$ gives its unique maximizer $\alpha_2$.
\end{proof}

\subsection{Exact node margins and root centering}

We now work at a maximizer. The preceding result removes the upper
and lower faces of the ordered simplex. There are multipliers
\begin{equation}
 \lambda_0=\lambda_m=0,\quad \lambda_j\ge0,\quad
 \partial_jJ=\lambda_{j-1}-\lambda_j,\quad
 \lambda_j(s_j-s_{j+1})=0\quad(j<m).
 \label{hier:eq:KKT}
\end{equation}
Write $m_j=\lambda_j-\lambda_{j-1}$. A maximal platform is a
maximal interval $[a,b]$ on which the products are equal. Its exterior
multipliers satisfy $\lambda_{a-1}=\lambda_b=0$.

\begin{lemma}[Node equations and Euler root identity]\label{hier:lem:KKT}
For every $j$,
\begin{equation}
 A_j\sum_{i\le j}L_iZ_i\prod_{q=i}^{j-1}\beta_q=p_jV_j+m_j.
 \label{hier:eq:nodebalance}
\end{equation}
Let $r_j=s_j/s_1$, $a_j=\log S_j$ and
$\mu_j=L_jZ_j$. Then
\begin{equation}
 T_1:=\sum_jr_jp_jV_j=\sum_j c_j\mu_j,
 \qquad c_j=r_ja_j. \label{hier:eq:rootcenter}
\end{equation}
Moreover $0<c_-(H)\le c_j,Z_j\le c_+(H)<\infty$.
\end{lemma}
\begin{proof}
Differentiation of the recursion gives
$\partial G_j/\partial s_j=A_j$ and
$\partial G_i/\partial G_j=\prod_{q=i}^{j-1}\beta_q$ for $i<j$.
Substitution into $\partial_jJ$ and \eqref{hier:eq:KKT} proves
\eqref{hier:eq:nodebalance}.
Every $G_i$ is homogeneous of degree one in the products.
Consequently
\[
 \sum_js_j\partial_jJ=\sum_js_jp_jV_j-\sum_iL_iZ_iG_i.
\]
The left side is zero because
$\sum_js_jm_j=\sum_{j<m}\lambda_j(s_j-s_{j+1})=0$.
Division by $s_1$ proves \eqref{hier:eq:rootcenter}. Finally
$2\le S_j\le H$ and \eqref{hier:eq:compact} bound $r_j$ above and below,
which gives the stated constants.
\end{proof}

\begin{proposition}[Poisson entropy at the saddle]\label{fp:prop:poisson}
For positive original lengths, let $s=s(L)$ be the ordered optimizer
in \eqref{profile:saddle}. With $\mathcal Q(z)=z\log z-z+1$,
\begin{equation}\label{fp:eq:poisson}
 J(L;s)=\sum_{i=1}^k L_i\mathcal Q(Z_i).
\end{equation}
Equivalently, $J$ is the relative entropy of independent Poisson counts
with means $L_iZ_i$ with respect to independent counts with means $L_i$.
\end{proposition}
\begin{proof}
The upper and lower simplex faces are inactive, so dilation of all
products is feasible in both directions, even on an equality face.
Homogeneity of the $G_i$ and stationarity give
\[
 0=\sum_j s_j\partial_{s_j}J
   =\sum_j s_jV_j-\sum_iL_iZ_iG_i.
\]
Substitute $G_i=\log Z_i$ into \eqref{profile:saddle}.
For a single count,
$\KL(\Pois(Lz)\Vert\Pois(L))=L\mathcal Q(z)$;
relative entropy adds over independent coordinates.
\end{proof}
The identity describes the physical tilt that supplies the exponential
factor of the profile. The persistence factors describe a second
selection, on the shorter time scale of the conditional bridges;
Section~\ref{fp:sec:interpretation} identifies its entropy cost.

\subsection{Complete unit-step chains: sign and effective time}

In this section only, assume
\[
 H\ge3,\qquad m=H-1,\qquad h_j=H-j+1,
 \qquad p_j=1\quad(1\le j\le m).
\]
The \emph{natural native node drift} is
\begin{equation}
 e_j=Z_jA_j-1. \label{hier:eq:nodedrift}
\end{equation}
The rank here is one. It must not be replaced by the larger rank of
an entire equal-product platform with its child left unsplit.

\begin{lemma}[A secant threshold with one valley]\label{hier:lem:valley}
For $x\ge2$ put
\[
 D(x)=x\log x-(x-1)\log(x-1),\qquad
 t(x)=1-\frac{\log D(x)}{\log x}.
\]
For $2\le x\le y$,
\begin{equation}
 t(x)\le\max\{\alpha_2,t(y)\}, \label{hier:eq:valley}
\end{equation}
strictly when $2<x<y$. More quantitatively, if $y\le H$,
$x\ge2+\eta$ and $y\ge x+\eta$, then
\begin{equation}
 \max\{\alpha_2,t(y)\}-t(x)\ge
 \Delta_H(\eta):=
 \frac{\eta^2\log(2\log2)}{2H^2(1+\log H)(\log H)^2}.
 \label{hier:eq:valleygap}
\end{equation}
\end{lemma}
\begin{proof}
Set $l(x)=\log(x/(x-1))$. Direct differentiation yields
\begin{align}
 t'(x)&=\frac{F(x)}{xD(x)(\log x)^2},\\
 F(x)&=D(x)\log D(x)-x\log x\,l(x),\\
 F'(x)&=l(x)\log D(x)
       +\log x\bigl(1/(x-1)-l(x)\bigr)>0.
 \label{hier:eq:Fprime}
\end{align}
Here $D(x)\ge2\log2>1$, and $\log(1+u)<u$ proves the second
term positive. Thus $t$ first decreases and then increases, allowing
either monotone case, with no flat interval. Its maximum on $[2,y]$
is attained at an endpoint and $t(2)=\alpha_2$. This proves
\eqref{hier:eq:valley} and its strict form.

For the quantitative assertion, throughout $[2,H]$,
\[
 F'(v)\ge\frac{\log(2\log2)}H,
 \qquad vD(v)(\log v)^2\le H(1+\log H)(\log H)^2.
\]
The latter bound follows from
$D(v)=\int_{v-1}^v(1+\log w)\,dw\le1+\log H$.
If $F(x)\ge0$, integrate $t'$ over $[x,x+\eta]$, using the
linear lower bound for $F(v)$ based at $x$, and then use monotonicity
up to $y$. If $F(x)\le0$, integrate backwards over $[x-\eta,x]$;
$F$ is nonpositive on $[2,x]$, so $t(2)\ge t(x-\eta)$.
Both integrations give exactly \eqref{hier:eq:valleygap}.
\end{proof}

\begin{theorem}[No negative native node drift]\label{hier:thm:positive}
On the complete unit-step chain, every positive length vector has
\begin{equation}
 e_j\ge0\quad(1\le j\le m),\qquad e_m>0.
 \label{hier:eq:allpositive}
\end{equation}
If $a$ is the first index of a maximal platform, then $e_a>0$ when
$a>1$. In addition, for every nonterminal index that is not first
in its platform,
\begin{equation}
 e_j\ge c_H:=2^{\Delta_H(\sqrt2-1)}-1>0.
 \label{hier:eq:uniformpositive}
\end{equation}
There is no assumed lower bound on length ratios, gaps between
platform products, or nonzero multipliers.
\end{theorem}
\begin{proof}
Let $[a,b]$ be a maximal platform with product $\sigma$. At its
first node $m_a=\lambda_a\ge0$. With
$I_{\rm in}=\sum_{i<a}L_iZ_i\prod_{q=i}^{a-1}\beta_q$, the exact
node equation and transport give
\begin{equation}
 \lambda_a=A_a(L_aZ_a+I_{\rm in})-V_a
 \le (Z_aA_a-1)V_a=e_aV_a. \label{hier:eq:firstsign}
\end{equation}
This proves $e_a\ge0$. If $a>1$, strict transport makes the
inequality strict, so $e_a>0$.

For $b<m$ put $C=h_{b+1}$ and $z=\exp(G_{b+1}/\sigma)$;
for $b=m$ put $C=z=1$. Within the platform
$S_i=h_i-C+z$ decreases by one at each step. At every node its
immediate child weight is $S_i-1$, including the last node of the
platform. Therefore
\begin{equation}
 A_i=D(S_i)/S_i,\qquad
 e_i=S_i^{\sigma-1}D(S_i)-1
     =\exp\bigl((\sigma-t(S_i))\log S_i\bigr)-1.
 \label{hier:eq:nativesecant}
\end{equation}
The first-node sign gives $\sigma\ge t(S_a)$; global compactness
gives $\sigma>\alpha_2$. Lemma~\ref{hier:lem:valley} propagates
$\sigma\ge t(S_i)$ through the platform. At the terminal node,
$e_m=2^{s_m}\log2-1>0$.

For $i>a$ with $i<m$, we have $S_a\ge S_i+1$. If $b=m$, then
$S_i=h_i\ge3$. Otherwise $u_b=s_{b+1}/\sigma>1/2$ and
$S_{b+1}\ge2$, whence $z>\sqrt2$ and $S_i\ge1+z>1+\sqrt2$.
Thus $S_i\ge2+(\sqrt2-1)$ in both cases. Apply
\eqref{hier:eq:valleygap} with $\eta=\sqrt2-1$ and then
\eqref{hier:eq:nativesecant}, using $S_i\ge2$, to obtain
\eqref{hier:eq:uniformpositive}.
\end{proof}

\begin{theorem}[Exact first-node drift scale]\label{hier:thm:effectivetime}
For the first index $a$ of any maximal platform on the complete
chain,
\begin{equation}
 e_a\asympH\frac{\lambda_a+V_{a-1}}{V_a}.
 \label{hier:eq:driftscale}
\end{equation}
For $a=1$ this is the exact equality $e_1=\lambda_1/L_1$.
When $a>1$, set
\[
 q_-=H^{-H},\qquad q_+=(1-1/H)^{\epsilon_H}<1.
\]
Then \eqref{hier:eq:driftscale} holds with explicit lower and upper
comparison constants $1-q_+$ and $q_-^{-1}$.
Consequently, with the infinite-quotient convention,
\begin{equation}
 1+\min\{L_a,e_a^{-2}\}
 \asympH
 1+\min\left\{L_a,
       \left(\frac{V_a}{\lambda_a+V_{a-1}}\right)^2\right\}.
 \label{hier:eq:efftime}
\end{equation}
For a non-first, nonterminal node, the left effective time is at
most $1+c_H^{-2}$.
\end{theorem}
\begin{proof}
On the unit-step chain $1/H\le\beta_q\le1-1/H$. Also
$\epsilon_H<1-s_q<1$. Formula \eqref{hier:eq:transport} implies, when
$a>1$,
\[
 q_-V_{a-1}\le I_{\rm in}/Z_a\le q_+V_{a-1}.
\]
There are fewer than $H$ factors in each transported term; the
upper bound uses at least one factor. Write
$w=I_{\rm in}/(Z_aV_{a-1})\in[q_-,q_+]$. The exact equality in
\eqref{hier:eq:firstsign}, before using transport, is
\begin{align}
 \lambda_a&=e_a(L_a+wV_{a-1})-(1-w)V_{a-1},\\
 e_a&=\frac{\lambda_a+(1-w)V_{a-1}}{L_a+wV_{a-1}}.
 \label{hier:eq:exactfirst}
\end{align}
Its numerator lies between
$(1-q_+)(\lambda_a+V_{a-1})$ and $\lambda_a+V_{a-1}$;
its denominator lies between $q_-V_a$ and $V_a$. This proves the
stated comparison. At $a=1$ there is no incoming intensity and
\eqref{hier:eq:nodebalance} gives the exact equality directly.
Taking inverse squares and truncating by $L_a$ preserves comparison
by fixed constants, including the zero-drift case. The final assertion
follows from \eqref{hier:eq:uniformpositive}.
\end{proof}

\begin{remark}
The effective time in \eqref{hier:eq:efftime} identifies a possible long
one-ended persistence scale. It does not evaluate reverse reserves,
terminal bands, or conditional persistence with a private incoming
state. In particular, it is not a claim that native packets can be
multiplied independently.
\end{remark}

\subsection{A hierarchy of maximal platform lengths}

We return to an arbitrary sparse chain $H=h_1>\cdots>h_m\ge2$.
No sign assertion for its individual native node drifts is needed.
For a maximal equal-product platform $[a,b]$, let its common product
be $\sigma$. Define
\begin{align}
 (C,z)&=(h_{b+1},S_{b+1}^{u_b}) &&(b<m),\\
 (C,z)&=(1,1) &&(b=m),\\
 n_i&=h_i-C\quad(a\le i\le b),& n&=n_a,& S&=S_a=n+z,\\
 A_*&=E_n(z)/S,&
 \mathfrak e_i&=(n_i+z)^{\sigma-1}E_{n_i}(z)-n_i.
 \label{hier:eq:platformparameters}
\end{align}
The symbol $\mathfrak e_i$ is a \emph{whole-platform} drift: it
uses the entire exterior rank $n_i$, with the child intact. It is
different from $e_i$ in \eqref{hier:eq:nodedrift}.

\begin{lemma}[Exact platform balance]\label{hier:lem:platformbalance}
With $I_{\rm in}=\sum_{i<a}L_iZ_i\prod_{q=i}^{a-1}\beta_q$,
\begin{equation}
 \sum_{i=a}^bL_i\mathfrak e_i
       =nV_{a-1}-A_*I_{\rm in}. \label{hier:eq:platformbalance}
\end{equation}
If $a>1$, then
\begin{equation}
 I_{\rm in}
 =\frac{\beta_{a-1}}{A_{a-1}}
    \bigl(p_{a-1}V_{a-1}-\lambda_{a-2}\bigr)
 \le\frac{\beta_{a-1}p_{a-1}}{A_{a-1}}V_{a-1}.
 \label{hier:eq:incoming}
\end{equation}
\end{lemma}
\begin{proof}
The sum of budget derivatives over $[a,b]$ is
$nV_{a-1}+\sum_{i=a}^bn_iL_i$. On the platform,
$G_i=\sigma\log(n_i+z)$ and its derivative is
$E_{n_i}(z)/(n_i+z)$, including the variation in $z$. The derivative
for an older block is $A_*\prod_{q=i}^{a-1}\beta_q$.
The sum of the KKT derivatives on the platform is
$\lambda_{a-1}-\lambda_b=0$. Substituting these derivatives proves
\eqref{hier:eq:platformbalance}.

At the single node $a-1$, the multiplier margin is
$m_{a-1}=\lambda_{a-1}-\lambda_{a-2}=-\lambda_{a-2}\le0$.
Equation \eqref{hier:eq:nodebalance}, followed by its routing factor
$\beta_{a-1}$, is exactly \eqref{hier:eq:incoming}. This step retains
the actual mixture of all older incoming intensities.
\end{proof}

\begin{lemma}[Whole-platform entropy surplus]\label{hier:lem:surplus}
If $a>1$, then
\begin{equation}
 n-A_*\frac{\beta_{a-1}p_{a-1}}{A_{a-1}}
 \ge\frac{np_{a-1}}{2H^2(1+\log H)^2}>0.
 \label{hier:eq:surplus}
\end{equation}
This estimate is independent of the size of the strict product gap
before the platform.
\end{lemma}
\begin{proof}
For real $r>0$ and $x\ge1$ let $D(r,x)=E_r(x)/r$. Integrating
$1+\log(x+t)$ gives
\begin{align}
 D(r,x)&=1+\log x+\frac1r\int_0^r\log(1+t/x)\,dt,\\
 1+\log x+\frac{r}{2(x+r)}&\le D(r,x)\le1+\log(x+r).
 \label{hier:eq:Dincrements}
\end{align}
For the lower bound use
$\log(1+t/x)\ge t/(x+t)\ge t/(x+r)$.
Write $p=p_{a-1}$, $u=u_{a-1}$ and $w=S^u$. Then
$1\le w\le S$, $p+w\le H$, and
\[
 \frac{\beta_{a-1}p}{A_{a-1}}=\frac{w}{D(p,w)},
 \qquad \frac n{A_*}=\frac S{D(n,z)}\ge\frac S{1+\log S}.
\]
Since $x/(1+\log x)$ is nondecreasing for $x\ge1$,
\begin{align}
 \frac n{A_*}-\frac{\beta_{a-1}p}{A_{a-1}}
 &\ge\frac w{1+\log w}-\frac w{D(p,w)}\notag\\
 &=\frac{w\,[D(p,w)-(1+\log w)]}
          {(1+\log w)D(p,w)}
 \ge\frac p{2H(1+\log H)^2}. \label{hier:eq:surplusstep}
\end{align}
Finally \eqref{hier:eq:Dincrements} implies $E_n(z)\ge n$, so
$A_*\ge n/H$. Multiplication of \eqref{hier:eq:surplusstep} by $A_*$
proves \eqref{hier:eq:surplus}. Strict separation of the products was used
only for the zero exterior multiplier, not to obtain its numerical
constant.
\end{proof}

\begin{theorem}[Maximal-platform length hierarchy]\label{hier:thm:hierarchy}
At every global ordered KKT maximizer on an arbitrary sparse chain,
each maximal platform $[a,b]$ satisfies
\begin{equation}
 V_{a-1}\le C_H\sum_{i=a}^bL_i,
 \qquad C_H=2H^3(1+\log H)^3.
 \label{hier:eq:hierarchy}
\end{equation}
For $a>1$, its signed incoming whole-platform drift satisfies
\begin{equation}
 A_*I_{\rm in}-nV_{a-1}
 \le-\frac{np_{a-1}}{2H^2(1+\log H)^2}\,V_{a-1}.
 \label{hier:eq:incominggap}
\end{equation}
The estimates require no positive length-ratio or product-gap bound.
\end{theorem}
\begin{proof}
Combining \eqref{hier:eq:incoming} and \eqref{hier:eq:surplus} gives
\eqref{hier:eq:incominggap}, and then \eqref{hier:eq:platformbalance} gives
\[
 \sum_{i=a}^bL_i\mathfrak e_i
 \ge\frac{np_{a-1}}{2H^2(1+\log H)^2}V_{a-1}.
\]
For every summand, $1\le n_i+z\le H$, $0<\sigma<1$ and
\eqref{hier:eq:Dincrements} give
$|\mathfrak e_i|\le H(1+\log H)$. Hence
\[
 V_{a-1}\le
 \frac{2H^3(1+\log H)^3}{np_{a-1}}\sum_{i=a}^bL_i
 \le C_H\sum_{i=a}^bL_i.
\]
For $a=1$, $V_0=0$, so \eqref{hier:eq:hierarchy} is immediate.
\end{proof}

\begin{corollary}[A long block in the final platform]\label{hier:cor:longblock}
Let $[a,m]$ be the final maximal platform, and select a largest block
$d$ within it. Then
\begin{equation}
 L_d\ge\frac{V_m}{(H-1)(1+C_H)}. \label{hier:eq:longblock}
\end{equation}
\end{corollary}
\begin{proof}
Set $W=\sum_{i=a}^mL_i$. Theorem~\ref{hier:thm:hierarchy} gives
$V_m\le(1+C_H)W$. The platform has at most $H-1$ blocks, so
$L_d\ge W/(H-1)$.
\end{proof}

\begin{remark}
The hierarchy is for whole maximal equal-product platforms. It does
not imply $V_{j-1}\ll_H L_j$ at an internal node of a platform.
Nor does \eqref{hier:eq:incominggap} make several native blocks within
that platform independent. Opposite whole-platform drift signs can
still cancel, and incoming retained words still carry private data.
\end{remark}

\subsection{An unread uncolored root reservoir}

This section is an elementary conditional local-mass statement for
the natural Poisson model at the current KKT optimizer. It is not a
path-survival theorem. In block $i$ take an independent homogeneous
Poisson point process of intensity $Z_i$ over length $L_i$, giving
$N_i\sim\Pois(\mu_i)$ with $\mu_i=Z_iL_i$. Any independent private
marks can be attached to these points. Let
\[
 V=V_m,\qquad X_1=\sum_i c_iN_i-T_1,
 \qquad c_i=r_i\log S_i.
\]
The centering $T_1=\sum_ic_i\mu_i$ is exact by
Lemma~\ref{hier:lem:KKT}.

Choose deterministically a block $d$ satisfying $L_d\ge\kappa_H V$
for some fixed $\kappa_H>0$. Either a globally largest block or the
block in Corollary~\ref{hier:cor:longblock} does so. In block $d$, choose
a deterministic subinterval of length $\tau L_d$, where
$0<\tau_0\le\tau\le\tau_1<1$ are fixed. Denote its uncolored count
by $N_R$ and its mean by $\mu_R=\tau\mu_d$.

Let $\mathcal O$ be the sigma-field of all points and private marks
\emph{outside} that interval. Then $N_R$ is independent of
$\mathcal O$. Write $N_i^{\rm out}=N_i$ for $i\ne d$ and
$N_d^{\rm out}=N_d-N_R$, with corresponding means
$\mu_i^{\rm out}=\mu_i$ and $\mu_d^{\rm out}=(1-\tau)\mu_d$.
The target $x$ must be measurable in the uncolored part of
$\mathcal O$. Assume, for fixed $A,B<\infty$,
\begin{equation}
 |x|\le A\log(2+V),\qquad
 \mathcal C_B=\bigcap_i
 \{|N_i^{\rm out}-\mu_i^{\rm out}|
                         \le B\sqrt{1+L_i}\}.
 \label{hier:eq:outsidecentral}
\end{equation}
The reservoir and its count have not been read by the target or by
any previously selected outside cost.

\begin{theorem}[Uniform conditional root atom]\label{hier:thm:reservoir}
Define the $\mathcal O$-measurable integer
\begin{equation}
 n_*(\mathcal O)=
 \left\lfloor\frac{T_1+x-\sum_i c_iN_i^{\rm out}}{c_d}\right\rfloor.
 \label{hier:eq:rootinteger}
\end{equation}
For $L_d\ge L_0(H,A,B,\kappa_H,\tau_0,\tau_1)$, on
$\mathcal C_B$ it is nonnegative and
\begin{align}
 |n_*-\mu_R|&\le C\sqrt{L_d},\label{hier:eq:rootcentral}\\
 \Prob(N_R=n_*\mid\mathcal O)&\asymp (1+V)^{-1/2}.
 \label{hier:eq:rootatom}
\end{align}
The constants depend only on the displayed fixed parameters. On
$N_R=n_*$,
\begin{equation}
 x-c_d<X_1\le x. \label{hier:eq:rootinterval}
\end{equation}
For every nonnegative $\mathcal O$-measurable random variable $F$,
including arbitrarily rare outside private-word weights,
\begin{equation}
 \E\bigl[F\ind_{\mathcal C_B}\ind_{\{N_R=n_*\}}\bigr]
 \asymp (1+V)^{-1/2}\E\bigl[F\ind_{\mathcal C_B}\bigr].
 \label{hier:eq:weightedatom}
\end{equation}
\end{theorem}
\begin{proof}
Subtract $c_d\mu_R$ inside \eqref{hier:eq:rootinteger}. Exact centering
then gives
\[
 n_*-\mu_R=
 \frac{x-\sum_ic_i(N_i^{\rm out}-\mu_i^{\rm out})}{c_d}
 +\varepsilon,\qquad -1<\varepsilon\le0.
\]
There are at most $H-1$ blocks. On $\mathcal C_B$,
\[
 \sum_i\sqrt{1+L_i}\le C_H\sqrt{1+V}
 \le C_{H,\kappa_H}\sqrt{L_d},
 \qquad \log(2+V)\le C_{H,\kappa_H}\sqrt{L_d}
\]
when $L_d\ge1$. Uniform bounds for $c_i$ prove
\eqref{hier:eq:rootcentral}. Since $\mu_R\ge c_H\tau_0L_d$, taking
$L_0$ sufficiently large makes $n_*$ positive.

For completeness, if $n=\mu_R+v$ with $|v|\le C\sqrt{\mu_R}$,
then
\[
 n\log(n/\mu_R)-n+\mu_R=O_C(1),
 \qquad n!\asymp \sqrt n\,(n/e)^n.
\]
Uniform Stirling inequalities therefore put
$e^{-\mu_R}\mu_R^n/n!$ between fixed positive multiples of
$\mu_R^{-1/2}$. Because $\mu_R\asymp L_d\asymp V$, this proves
\eqref{hier:eq:rootatom} pointwise on every admitted outside fiber.
The floor in \eqref{hier:eq:rootinteger} gives \eqref{hier:eq:rootinterval}
exactly. Finally condition on $\mathcal O$ and use the pointwise
conditional estimate to obtain \eqref{hier:eq:weightedatom}.
\end{proof}

\begin{proposition}[Compatibility with exact margin-carry bands]
Assume $m\ge2$. Suppose, in addition to the hypotheses of
Theorem~\ref{hier:thm:reservoir}, that a subsequent retained word satisfies the exact recursions
\begin{equation}
 X_{j+1}=(X_j-Y_j)/u_j\quad(j<m)
 \label{hier:eq:rootrecursion}
\end{equation}
and the bands are represented exactly as
\begin{equation}
 Y_j=m_j+y_j+\xi_j-\xi_{j-1}/u_{j-1},\quad j<m,
 \qquad \xi_0=0,
 \label{hier:eq:carrybands}
\end{equation}
where the initial term with $u_0$ is absent. Choose the uncolored
root target $x=\sum_{j<m}r_jy_j$. Then
\begin{equation}
 X_m=\frac{X_1-x}{r_m}-\lambda_{m-1}
                         -\frac{\xi_{m-1}}{u_{m-1}}.
 \label{hier:eq:carrylast}
\end{equation}
The identity is unaffected by which block supplies the root count.
\end{proposition}
\begin{proof}
Since $r_{j+1}=r_ju_j$, recursion \eqref{hier:eq:rootrecursion} gives
$r_mX_m=X_1-\sum_{j<m}r_jY_j$. Complementarity in
\eqref{hier:eq:KKT} gives $\sum_{j\le m}r_jm_j=0$ and
$m_m=-\lambda_{m-1}$, hence
$\sum_{j<m}r_jm_j=r_m\lambda_{m-1}$. The carry terms telescope:
\[
 \sum_{j<m}r_j(\xi_j-\xi_{j-1}/u_{j-1})
       =r_{m-1}\xi_{m-1}.
\]
Substitute these equalities and $x=\sum_{j<m}r_jy_j$ to obtain
\eqref{hier:eq:carrylast}. The argument contains no assumption about
the position of the selected reservoir.
\end{proof}

\begin{remark}[The additional conditional obligation]
After conditioning on $N_R=n_*$, the reservoir points are uniform
order statistics on the fixed interval and their private marks
retain their original law. A later packet $P$ that reads that
reservoir needs a further pointwise conditional estimate at the
prescribed count. For example, if on $\mathcal C_B$,
\[
 c\,G(\mathcal O)\le
 \E[P\mid\mathcal O,N_R=n_*(\mathcal O)]
 \le C\,G(\mathcal O),
\]
then one more conditioning step extends \eqref{hier:eq:weightedatom}
with $P$ on the left and $G$ on the right. A marginal scalar
estimate averaged over $N_R$ cannot substitute for this hypothesis.
Neither an ordinary bridge factor nor a fractional persistence
factor is part of Theorem~\ref{hier:thm:reservoir}.
\end{remark}

